\section{三角関数}
  その仕組みを確認するために,まず加法定理について復習しよう.
  \begin{align*}
    \sin(\alpha+\beta)
      &=\sin\alpha\cos\beta+\cos\alpha\sin\beta\\
    \cos(\alpha+\beta)
      &=\cos\alpha\cos\beta-\sin\alpha\sin\beta\\
    \tan(\alpha+\beta)
      &=\frac{\tan\alpha+\tan\beta}
              {1-\tan\alpha\tan\beta}
  \end{align*}
  ここで $\alpha=\beta=x$ とすれば,倍角公式
  \begin{align*}
    \sin 2x&=2\sin x\cos x\\
    \cos 2x&=2\cos^2x-1=1-2\sin^2x\\
    \tan 2x&=\frac{2\tan x}{1-\tan^2x}
  \end{align*}
  を得る.また,倍角公式から
  \begin{align*}
    \sin^2x&=\frac{1-\cos 2x}{2}\\
    \cos^2x&=\frac{1+\cos 2x}{2}\\
    \tan^2x&=\frac{1-\cos 2x}{1+\cos 2x}
  \end{align*}
  を導くことができる.また,加法定理から,\,$3x=2x+x$とすることで,
  \begin{align*}
    \sin 3x&=3\sin x-4\sin^3x\\
    \cos 3x&=4\cos^3x-3\cos x
  \end{align*}
  が求められる.